\documentclass[10pt,a4paper, ngerman]{beamer}

\usepackage[ngerman]{babel}
\usepackage[T1]{fontenc}
\usepackage[utf8]{inputenc}
\usepackage{varioref}
\usepackage{hyperref}
\usepackage{cleveref}
\usepackage{amsmath}
\usepackage{amsfonts}
\usepackage{amssymb}
\usepackage{makeidx}
\usepackage{graphicx}
\usepackage{csquotes}
\usepackage{listings}
\usepackage{color}
\usepackage{xcolor}
\usepackage[most]{tcolorbox}
\usepackage{amssymb}
\usepackage{lmodern}
\usepackage{verbatim}

% Umlaute und ß in Listings
\lstset{basicstyle=\ttfamily}
\lstset{literate=%
  {Ö}{{\"O}}1
  {Ä}{{\"A}}1
  {Ü}{{\"U}}1
  {ß}{{\ss}}1
  {ü}{{\"u}}1
  {ä}{{\"a}}1
  {ö}{{\"o}}1
}

% Farben für Listings
\definecolor{codegreen}{rgb}{0,0.6,0}
\definecolor{codegray}{rgb}{0.5,0.5,0.5}
\definecolor{codepurple}{rgb}{0.58,0,0.82}
\definecolor{backcolour}{rgb}{0.95,0.95,0.92}
 
\lstdefinestyle{mystyle}{
    backgroundcolor=\color{backcolour},   
    commentstyle=\color{codegreen},
    keywordstyle=\color{magenta},
    numberstyle=\tiny\color{codegray},
    stringstyle=\color{codepurple},
    basicstyle=\footnotesize,
    breakatwhitespace=false,         
    breaklines=true,                 
    captionpos=t,                    
    keepspaces=true,                 
    numbers=left,                    
    numbersep=5pt,                  
    showspaces=false,                
    showstringspaces=false,
    showtabs=false,                  
    tabsize=2,
	language=[LaTeX]{TeX}
}

\lstset{style=mystyle}

\definecolor{hbbkBlueish}{HTML}{5e7b87}
\definecolor{hbbkReal}{HTML}{1A4354}
\setbeamercolor{logo}{bg=white}  %controls the color of the logo area
\setbeamercolor{author in sidebar}{fg=white} 
\setbeamercolor{palette primary}{bg=hbbkBlueish,fg=white}
\setbeamercolor{palette secondary}{bg=hbbkBlueish,fg=white}
\setbeamercolor{palette tertiary}{bg=hbbkBlueish,fg=white}
\setbeamercolor{palette quaternary}{bg=hbbkBlueish,fg=white}
\setbeamercolor{structure}{fg=hbbkBlueish} % itemize, enumerate, etc
\setbeamercolor{section in toc}{fg=hbbkReal} % TOC sections
\setbeamertemplate{navigation symbols}{}%remove navigation symbols

\setbeamercovered{invisible}

\usetheme{PaloAlto}
\renewcommand{\footnotesize}{\small}
\newcommand{\pftn}[1]{\let\thefootnote\relax\footnotetext{\tiny #1}}
\newcommand{\ftn}[2]{\footnote[#1]{\tiny #2}}
\newenvironment{hlbox}{\begin{tcolorbox}[enhanced,colback=white,colframe=white,sharpish corners,fuzzy halo=0.5mm with lightgray]}{\end{tcolorbox}}

\logo{\includegraphics[width=0.155\linewidth]{hbbk-logo}}


\usepackage{eurosym}

\AtBeginSection{\frame{\frametitle{Gliederung}\tableofcontents[currentsection]}}

\author{Luca Kiebel}
\title{Anleihen}
%\subtitle{subtitle}
\date{\today}
\institute[HBBK]{Hans-Böckler-Berufskolleg}
\setlength{\itemsep}{10pt}
\begin{document}
\begin{frame}
\titlepage
\end{frame}

\section{Einleitung}
\subsection{Was ist eine Anleihe?}
\begin{frame}{Was ist eine Anleihe?}{Einleitung}
\pftn{https://de.wikipedia.org/wiki/Anleihe}
\begin{itemize}
	\item Wertpapier
	\item Gibt Gläubiger Recht auf Rückzahlung + Zinsen
	\item Helfen Schuldner zur Fremdfinanzierung
	\item Helfen Gläubiger zur Kapitalanlage
\end{itemize}
\end{frame}

\subsection{Arten von Anleihen}
\begin{frame}{Arten von Anleihen}{Einleitung}
\pftn{https://de.wikipedia.org/wiki/Anleihe}
\begin{itemize}
	\item Staatsanleihen
	\begin{itemize}
		\item<2-> Staat ist Schuldner
		\item<2-> Bonität bestimmt durch Staatshaushalt
	\end{itemize}
	\item Unternehmensanleihen
	\begin{itemize}
		\item<3-> Unternehmen ist Schuldner
		\item<3-> Bonität bestimmt durch Rating
	\end{itemize}
	\item Pfandbriefe
	\begin{itemize}
		\item<4-> Hypotheken auf Immobilien
		\item<4-> Bonität ist durch Pfand bestimmt
	\end{itemize}
\end{itemize}
\end{frame}

\section{Rendite}
\begin{frame}{Rendite}{Berechnung}
\[Rendite = 100*\left ( \frac{\left ( Zinskoupon + \frac{Verkaufskurs - Kaufkurs}{Laufzeit} \right )}{Kaufkurs} \right )\]
\pftn{https://www.gevestor.de/details/effektivverzinsung-bei-anleihen-so-wird-sie-berechnet-619805.html}
\pause
\begin{block}{Beispiel: Bundesanleihe (10 Jahre)}
\textbf{Zinskoupon}: 0,5 \% \textbf{Kaufkurs}: 65 \% \textbf{Verkaufskurs}: 75 \% \\
--> \textbf{Rendite}: 2,31 \%
\end{block}
\pftn{https://www.deutsche-finanzagentur.de/de/institutionelle-investoren/bundeswertpapiere/bundesanleihen/}
\pftn{https://www.finanzen.net/zinsen/10j-Bundesanleihen}
\end{frame}

\section{Liquidität}
\begin{frame}{Liquidität}
\begin{itemize}
	\item Können jederzeit an der Börse gehandelt werden
	\item Teilweise Handel von Person zu Person möglich \pause
	\item Kosten für Handel an der Börse
\end{itemize}
\pftn{https://www.finanztip.de/anleihen/staatsanleihen/}
\end{frame}

%\section{Kosten}
%\begin{frame}
%
%\end{frame}

\section{Voraussetzungen}
\subsection{Mindesteinlage}
\begin{frame}{Mindesteinlage}{Vorraussetzungen}
\pftn{Judith Engst \& Janne Kipp{,} \textit{Erfolgreiche Geldanlage für Dummies}{,} Wiley, 2017 (Seite 156)}
\begin{itemize}
	\item Deutsche Bundesanleihe: 1 cent\ftn{1}{https://www.deutsche-finanzagentur.de/de/factsheet/sheet-detail/productdata/sheet/DE0001102440/}
	\item Unternehmensanleihen: Meist höher (>20.000 \euro)
\end{itemize}
\end{frame}


\section{Nachhaltigkeit}
\begin{frame}{Nachhatigkeit}
\begin{itemize}
	\item Kommt auf den Schuldner an
	\item Green Bonds: Umweltfreundlich
\end{itemize}
\pftn{https://www.thebalance.com/what-are-green-bonds-417154}
\pftn{Judith Engst \& Janne Kipp{,} \textit{Erfolgreiche Geldanlage für Dummies}{,} Wiley, 2017 (Seite 180)}
\end{frame}

\end{document}
